\subsection{\texorpdfstring{E$_8$}{E8} as downstream scale grammar}

In the present paper E$_8$ is not used as a primitive cause. The retained parameter is
\[
\kappa_{E8}:=\frac{\gamma}{\ln(248/60)}=\frac{5/6}{\ln(248/60)}.
\]
On the retained branch this means
\[
\begin{aligned}
\kappa_{E8}&\approx \num{0.58723},
&
f_a&\approx \num{8.86e10}\,\mathrm{GeV},
\\
m_a&\approx \num{65.19}\,\mu\mathrm{eV},
&
\nu_a&\approx \num{15.764}\,\mathrm{GHz}.
\end{aligned}
\]
The intended reading is a downstream scale grammar:
\begin{itemize}[leftmargin=1.5em, topsep=3pt, itemsep=2pt]
    \item the carrier packet fixes $\gamma=5/6$,
    \item $\kappa_{E8}$ then defines an optional ladder parameter,
    \item the ladder is allowed to organize scale spacing,
    \item but it does not replace the carrier theorem or the seam data.
\end{itemize}
In the fuller appendix-level grammar one may package a downstream block scale as
\[
X(n,r,I_1)=\frac{\Mplunred}{8}\,\sin^2\theta_{13}
\left(\frac{60-2n}{60}\right)^{\kappa_{E8}}
e^{-(8-r)/64}e^{-12\pi I_1},
\]
where $n$ is the ladder stage, $r$ a mild block-multiplicity label, and $I_1$ the dominant block trace. This makes the reading discipline explicit: the stage index orders the ladder, the trace controls most of the exponential suppression, and the extra block label only refines the scale. The denominator $64$ is part of that downstream grammar and is not interpreted here as a Hilbert-space dimension, state count, or hidden spinor multiplicity. E$_8$ is therefore a downstream scale grammar, not a universal mass formula.

\subsection{Full stage atlas}

The full real-domain atlas up to $n=29$ extends the main-text robust table with restricted-scan candidates and IR tail conjectures. Rows are classified as \emph{robust} (restrictive grammar, unambiguous target), \emph{restricted-scan} (plausible but less constrained), or \emph{tail conjecture} (new trace class $I_1\in\{\frac{11}{12},\frac{7}{6}\}$ not yet derived from block structure).

\begin{table}[H]
\centering
\scriptsize
\setlength{\tabcolsep}{2pt}
\begin{tabularx}{\linewidth}{@{}>{\centering\arraybackslash}p{0.035\linewidth}>{\centering\arraybackslash}p{0.035\linewidth}>{\raggedright\arraybackslash}p{0.065\linewidth}>{\raggedleft\arraybackslash}p{0.12\linewidth}>{\raggedleft\arraybackslash}p{0.12\linewidth}>{\raggedright\arraybackslash}p{0.09\linewidth}>{\raggedright\arraybackslash}p{0.17\linewidth}Y@{}}
\toprule
$n$ & $r$ & $I_1$ & $X$ (GeV) & Target & Residual & Identification & Class \\
\midrule
$3$ & $8$ & $1/2$ & $\num{2.159e8}$ & $\num{2.160e8}$ & $0.06\%$ & $\mu_{\cthree}$ RG crossing & robust \\
$5$ & $5$ & $1/12$ & $\num{1.307e15}$ & $\num{1.31e15}$ & $0.2\%$ & seesaw $M_R$ & robust \\
$6$ & $6$ & $1$ & $\num{1.272}$ & $\num{1.27}$ & $0.2\%$ & charm band & restricted-scan \\
$8$ & $1$ & $17/16$ & $\num{0.1059}$ & $\num{0.10566}$ & $0.2\%$ & muon band & restricted-scan \\
$10$ & $1$ & $1/3$ & $\num{8.826e10}$ & direct search & --- & PQ scale $f_a$ & robust \\
$12$ & $2$ & $41/48$ & $\num{246.35}$ & $\num{246.22}$ & $0.05\%$ & EW reference scale & robust \\
$14$ & $4$ & $17/16$ & $\num{0.0921}$ & $\num{0.0927}$ & $0.6\%$ & $f_\pi$ band & robust \\
$15$ & $4$ & $1$ & $\num{0.935}$ & $\num{0.938}$ & $0.3\%$ & hadron anchor & robust \\
$21$ & $4$ & $1/6$ & $\num{3.051e13}$ & $\num{3.06e13}$ & $0.3\%$ & inflation $M$ & robust \\
$21$ & $5$ & $41/48$ & $\num{171.8}$ & $\num{172.6}$ & $0.4\%$ & top quark & robust \\
$24$ & $7$ & $5/8$ & $\num{7.895e5}$ & $\num{7.920e5}$ & $0.3\%$ & $\mu_{\phiz}$ RG crossing & robust \\
$25$ & $5$ & $1$ & $\num{0.498}$ & $\num{0.498}$ & $0.0\%$ & kaon band & robust \\
$25$ & $7$ & $41/48$ & $\num{125.5}$ & $\num{125.25}$ & $0.2\%$ & Higgs mass & robust \\
$26$ & $2$ & $1$ & $\num{0.417}$ & $\num{0.42}$ & $0.7\%$ & $\sqrt{\sigma}$ confinement & robust \\
$27$ & $6$ & $41/48$ & $\num{91.6}$ & $\num{91.19}$ & $0.4\%$ & $Z$ boson & restricted-scan \\
$28$ & $1$ & $7/6$ & $\num{5.105e-4}$ & $\num{5.110e-4}$ & $0.1\%$ & electron & tail conjecture \\
$29$ & $1$ & $11/12$ & $\num{4.210}$ & $\num{4.18}$ & $0.7\%$ & bottom quark & tail conjecture \\
\bottomrule
\end{tabularx}
\caption{Full E$_8$ stage atlas on the canonical branch up to the real-domain boundary $n=29$.}
\label{tab:e8-full-atlas}
\end{table}

\subsection{Compression identities on the minimal branch}

The main-text compression identities are developed in the dedicated section on compression identities and carrier exclusion. This appendix subsection collects additional arithmetic continuations and the compression chain summary.

\begin{remark}[Abelian package from occupancy and carrier norm]
On the canonical branch,
\[
\Oadm\gamma=48\cdot\frac56=40.
\]
Thus the characteristic coefficient $41$ is compressed to
\[
41=\Oadm\gamma+1=40+1,
\]
that is, fermionic occupancy trace plus one light Higgs doublet contribution.
Writing $b_Y:=5\bone/3$ for the usual hypercharge-normalized coefficient and using the block-index relation $k_B=\frac32 I_1(B)$ to define $I_1^{\mathrm{EW}}:=\frac{5}{24}\bone$ and $k_{\mathrm{EW}}:=\frac{5}{16}\bone$, one obtains the unified chain
\begin{equation}
\label{eq:41-chain}
48\,I_1^{\mathrm{EW}}
=32\,k_{\mathrm{EW}}
=10\,\bone
=\Oadm\gamma+1
=41.
\end{equation}
Equivalently,
\[
I_1^{\mathrm{EW}}=\frac{41}{48},
\qquad
k_{\mathrm{EW}}=\frac{41}{32},
\qquad
\bone=\frac{41}{10},
\qquad
b_Y=\frac{41}{6}.
\]
This is an algebraic signature rather than an accidental coincidence: the same integer $41$ appears simultaneously as electroweak block trace, E$_8$ exponent, UV hypercharge trace, and carrier occupancy times norm plus one.
\end{remark}

\begin{remark}[Rank lock and single-seed compression: main-text reference]
The rank-lock identity $2^{g-3}=g-1\Rightarrow g=5$ is proved as \Cref{prop:rank-lock} in the main text.
The single-seed compression and the Cabibbo inversion $\phiz=\frac12(1-\sqrt{1-4\lambda_C^2})$ with all derived inter-observable relations are given in \Cref{eq:cabibbo-inversion,eq:lambda-c-master}. The carrier exclusion theorems for wrong discrete worlds are in \Cref{thm:carrier-exclusion}.
\end{remark}

\subsection{Optional arithmetic continuations}

The next relations mix the main-text carrier with optional appendix-level E$_8$, inflation, or geometric notation. They are useful compression patterns, but their interpretation should remain conditional or conjectural.

\begin{remark}[60, 48, and \texorpdfstring{$5/4$}{5/4} on the optional ladder]
If the optional scale ladder is started at $D_1=60$, then
\[
\frac{D_1}{\Oadm}=\frac{60}{48}=\frac54=B\gamma.
\]
On the minimal branch one may also write the same number as
\[
\frac54=\frac{\gcar}{\gcar-1}\qquad (\gcar=5).
\]
This compresses the E$_8$ start index, the admissible chiral count, the mirror factor, and the two-defect coefficient to the same rational factor.
\end{remark}

\begin{remark}[Factorization of the two defect coefficient and arithmetic completion of $48$ and $60$]\label{rem:defect-factorization-gcd-lcm}
From the main text identity
\[
\delta_2=\frac{5}{4}\,\deltatop^2,
\qquad
\deltatop=\Oadm\cthree^4,
\]
one obtains
\[
\delta_2=2880\,\cthree^8.
\]
On the canonical branch this factor admits the exact arithmetic decompositions
\[
2880=\Dstart\Oadm=60\cdot48,
\qquad
2880=|R(E_8)|\,\dim \mathfrak g_{\SM}=240\cdot12.
\]
The same pair $(\Oadm,\Dstart)=(48,60)$ also satisfies
\[
\gcd(\Oadm,\Dstart)
=
12
=
\dim \mathfrak g_{\SM},
\]
and
\[
\operatorname{lcm}(\Oadm,\Dstart)
=
\frac{\Oadm\Dstart}
{\gcd(\Oadm,\Dstart)}
=
240
=
|R(E_8)|.
\]
So the gauge algebra dimension is the common divisor of the carrier-side counting pair, while the E$_8$ root count is its least common multiple. The arithmetic identities are exact. Here the value $12$ is the Lie-algebra dimension
\[
\dim\mathfrak g_{\SM}=8+3+1,
\]
not a derivation from a discrete permutation group of order $12$; at most the $3+2$ split carries such a discrete shadow. Their E$_8$ reading remains appendix level and conditional in the sense fixed elsewhere in the manuscript.
\end{remark}

\begin{conjecture}[E$_8$ root-count reading]
On the same optional branch,
\[
|R(E_8)|=248-8=240=5\cdot48=5\,\Oadm.
\]
This suggests that the E$_8$ root count may be read as a fivefold unfolding of the admissible chiral packet, although that interpretation is not used as a main-text theorem.
\end{conjecture}

\begin{remark}[Reactor seed of the optional scale ladder]
If one writes the optional ladder as
\[
\varphi_n=\phiz e^{-\gamma}\left(\frac{D_n}{D_1}\right)^{\kappa_{E8}},
\qquad
D_n=60-2n,
\]
then the start value is precisely the reactor-angle seed,
\[
\varphi_n=\sin^2\theta_{13}\left(\frac{D_n}{D_1}\right)^{\kappa_{E8}}.
\]
Thus the appendix-level ladder may be read as a discrete continuation of the same seed that controls $\theta_{13}$.
\end{remark}

\begin{remark}[Conditional inflation-side compression]
If one supplements the appendix with the conditional Starobinsky-side relations
\[
A_s=\frac{N^2}{24\pi^2}\left(\frac{M}{\Mpl}\right)^2,
\qquad
\frac{M}{\Mpl}=\sqrt{8\pi}\,\cthree^4,
\]
then on the minimal branch
\[
A_s=\frac{N^2}{3\pi}\cthree^8=\frac{N^2}{8640\pi}\,\delta_2,
\]
and also
\[
\frac{M}{\Mpl}=\frac{\sqrt{8\pi}}{\Oadm}\,\deltatop.
\]
This is recorded only as an appendix-level compression identity because inflation is outside the present main-text claim surface.
\end{remark}

\begin{remark}[Dyonic intercept and the $\alpha$ kernel]
From the dyonic appendix the local astrophysical $\beta$ function reads
\[
\beta_{\mathrm{BH}}(r)=\frac{Q_eQ_m}{256\pi^4 r^2}=\frac{\deltatop}{3}\frac{Q_eQ_m}{r^2}=16\cthree^4\frac{Q_eQ_m}{r^2}.
\]
This is a cross-link between the EHT intercept channel and the $\alpha$ kernel: the same topological coefficient $\deltatop=48\cthree^4$ that pushes $\alpha$ into the precision zone also controls the local dyonic $\beta$ function. The local and cosmological sectors are therefore not two loose motifs but the same coefficient in two projections.
\end{remark}

\begin{remark}[Conditional occupancy route to gravity]
If the relative Einstein term is of the form
\[
S_{\mathrm{EH}}^{\mathrm{rel}}
=
\frac{\chigeo^2}{192\pi^2}\tr_{\mathcal H_{\mathrm{ch}}}\mathbf 1
\int d^4x\sqrt{-g}\,R,
\]
and if the family bridge identifies
\[
\tr_{\mathcal H_{\mathrm{ch}}}\mathbf 1=\Oadm=48,
\]
then
\[
\Mpl^2=\frac{\chigeozero^2}{2\pi^2}.
\]
This is the occupancy-to-gravity route alluded to in the main text.
\end{remark}

\begin{remark}[Compression chain]
The appendix-level compression pattern may therefore be summarized as
\[
Y \leftrightarrow \inv,
\qquad
\gamma=\frac56,
\qquad
B\gamma=\frac54,
\qquad
\frac{D_1}{\Oadm}=\frac54,
\qquad
|R(E_8)|=5\,\Oadm,
\]
\[
\bone=\frac{\Oadm\gamma+1}{10},
\qquad
\Omega_b+\beta_{\mathrm{rad}}=\phiz,
\qquad
A_s=\frac{N^2}{8640\pi}\,\delta_2
\quad
\text{(conditional inflation-side import)}.
\]
\end{remark}

\subsection{Conditional minimal-parameter picture}

The next package is stronger but also more conditional. It assumes that the carrier rank, family multiplicity, and seam normalization are all read as closed rather than only partially inherited.

\begin{principle}[Conditional minimal-parameter compression]
Assume the canonical branch
\[
\gcar:=\rank E=5,
\qquad
N_{\mathrm{fam}}=3,
\qquad
\cthree=\frac{1}{8\pi}.
\]
Then the numerical kernel may be rewritten almost entirely in terms of
\[
(\pi,\gcar,N_{\mathrm{fam}})=(\pi,5,3).
\]
Indeed, on this branch one has
\[
\gamma=\frac56=\frac{\gcar}{\gcar+1},
\qquad
\dim S^+=2^{\gcar-1},
\qquad
\Oadm=N_{\mathrm{fam}}2^{\gcar-1}=48,
\]
\[
\phibase=\frac{1-\gamma}{\pi}=\frac{1}{(\gcar+1)\pi},
\qquad
\deltatop=\Oadm\,\cthree^4,
\qquad
\delta_2=\Dstart\,\Oadm\,\cthree^8,
\]
with
\[
\Dstart=\gcar\,\dim\mathfrak g_{\SM}=5\cdot 12=60,
\qquad
\Oadm=(\gcar-1)\dim\mathfrak g_{\SM}=4\cdot 12=48.
\]
\end{principle}

\begin{remark}[Compressed dependency chain]
Under the same conditional closure, the appendix-level numerical chain reads
\[
(\widetilde M\to M,\Seam,\inv,\chiseed)
\Longrightarrow
E_3\oplus E_2
\Longrightarrow
Y \leftrightarrow \inv
\Longrightarrow
S(U(3)\times U(2))
\Longrightarrow
S^+,
\]
\[
\Longrightarrow
N_{\mathrm{fam}}=3
\Longrightarrow
\Oadm=48
\Longrightarrow
(\deltatop,\delta_2,\phiz)
\Longrightarrow
(\alpha,\lambda_C,\beta_{\mathrm{rad}},\Omega_b,\sin^2\theta_{13}),
\]
with the optional E$_8$, inflationary, and gravity-side sectors then read as downstream continuations of the same compressed core.
\end{remark}

\section{Relative APS and superconnection setup}
\label{app:aps}

The geometric completion uses APS-type relative data \cite{aps1975}. In practice this means:
\begin{itemize}[leftmargin=1.5em, topsep=3pt, itemsep=2pt]
    \item one works with a declared pair $(D,D_{\mathrm{ref}})$ rather than a single unanchored operator,
    \item boundary conditions on the seam $\Seam$ are part of the definition of the relative index,
    \item the $\eta$-term keeps track of the residual spectral asymmetry,
    \item the superconnection package is used in the same spirit as the spectral-action tradition \cite{chamseddineConnes1997}, but here always with an explicit reference subtraction.
\end{itemize}

This appendix now serves as the proof motor for the geometric reconstruction statements in the main text: its role is to make the common reference structure explicit enough that the reconstruction theorem and the local spectral-scale branch are not left hanging on a slogan.

\section{Relative reflection positivity with fermions}
\label{app:os-fermions}

\subsection{Normalized quotient form}

The relative Euclidean theory used in the main closure section is not interpreted as a difference of measures. Instead it is organized as the normalized quotient
\[
Z_{\mathrm{rel}}[J,\eta,\bar\eta]
=
\frac{Z_{\mathrm{adm}}[J,\eta,\bar\eta]}{Z_{\mathrm{ref}}[0,0,0]}.
\]
Because the denominator is source independent on the admissible sector, the reference theory acts only as a fixed normalizer and does not interfere with reflection positivity.

\subsection{Bosonic and fermionic positivity on \texorpdfstring{$\Padm$}{Padm}}

The commutation relations
\[
[\Theta,P_\Theta]=0,
\qquad
\Theta\Padm=\Padm\Theta
\]
ensure that reflection preserves the admissible subspace. The bosonic part then follows the standard Markov / reflection-stable Osterwalder--Schrader argument on admissible observables. For fermions, the admissible CAR kernel may be written as a reflected quadratic form; equivalently one may express the same condition as positivity of the associated Pfaffian structure on the admissible sector. Hence for every admissible observable $O$ supported in $\tau>0$,
\[
\langle \Theta O,O\rangle_{\mathrm{rel}}
=
\frac{1}{Z_{\mathrm{ref}}[0,0,0]}
\langle \Theta O,O\rangle_{\mathrm{adm}}
\ge 0.
\]
This appendix section is the mechanism behind \Cref{lem:admissible-reflection-positivity,thm:reflection-positivity-admissible,thm:lorentzian-reconstruction}; it is recorded here so that the main-text reconstruction theorem is not supported only by a one-line axiom call.

